
\subsubsection{5.1 RQ1: Layer Heterogeneity (H-E1: PASS)}

At primary epsilon = 0.01, the 32-layer sparsity profile measured from 512 Alpaca samples exhibits CV = 0.544, well above the gate threshold of 0.3 (mean sparsity = 0.0227, standard deviation = 0.0124). The Kendall's tau between Alpaca and WikiText-103 rankings is tau\_calibration = 0.786 (p = 2.03e-13), above the threshold of 0.6. The rank correlation between 128-token and 512-token Alpaca profiles is tau\_length = 0.899 (p = 2.93e-20).

A systematic depth gradient is observable in the raw per-layer data: layers 0--2 exhibit the highest mean sparsity values (approximately 0.055--0.082 at epsilon = 0.01), while layers 28--31 exhibit the lowest (approximately 0.011--0.014). Intermediate layers show a roughly monotone decline across depth with some local variation.

\begin{table*}[htbp]
\centering
\resizebox{\dimexpr 2\columnwidth+\columnsep\relax}{!}{%
\begin{tabular}{llllll}
\toprule
epsilon & CV (Alpaca, len=512) & tau\_calibration & p-value (tau\_cal) & tau\_length & p-value (tau\_len) \\
\midrule
0.001 & 0.549 & 0.790 & 1.33e-13 & 0.883 & 5.41e-19 \\
**0.010** & **0.544** & **0.786** & **2.03e-13** & **0.899** & **2.93e-20** \\
0.050 & 0.528 & 0.778 & 4.64e-13 & 0.875 & 2.08e-18 \\
0.100 & 0.484 & 0.782 & 3.08e-13 & 0.879 & 1.07e-18 \\
\bottomrule
\end{tabular}
}
\end{table*}

All four epsilon values pass both gate conditions. H-E1 gate result: \textbf{PASS}.

\begin{figure}[htbp]
\centering
\includegraphics[width=\columnwidth]{figures/sparsity_profile.png}
\caption{Sparsity profile across 32 layers}
\label{fig:sparsity_profile}
\end{figure}

\textit{Figure 1. Per-layer mean sparsity at epsilon = 0.01 for Alpaca (length 512) and WikiText-103 (length 512), across all 32 MLP gate\_proj layers of LLaMA-3.1-8B.}

\begin{figure}[htbp]
\centering
\includegraphics[width=\columnwidth]{figures/epsilon_sensitivity.png}
\caption{Epsilon sensitivity}
\label{fig:epsilon_sensitivity}
\end{figure}

\textit{Figure 2. CV and tau\_calibration across all four epsilon values. All values exceed gate thresholds (dashed lines).}

\subsubsection{5.2 RQ2: Cross-Distribution Stability (H-M1: PASS)}

Across four diverse calibration distributions measured at primary epsilon = 0.01, ICC(3,k) = 0.9846 (95\% CI: [0.97, 0.99]), far above the gate threshold of 0.75. All six pairwise Kendall's tau values exceed 0.6; the minimum is tau = 0.7339 for the WikiText-103 vs. SST-2 pair (p = 2.78e-11).

\begin{table}[htbp]
\centering
\resizebox{\columnwidth}{!}{%
\begin{tabular}{lll}
\toprule
Pair & Kendall's tau & p-value \\
\midrule
Alpaca vs. WikiText-103 & 0.7863 & 2.03e-13 \\
Alpaca vs. SST-2 & 0.9395 & 3.35e-24 \\
Alpaca vs. MNLI & 0.9476 & 3.51e-25 \\
WikiText-103 vs. SST-2 & **0.7339** & 2.78e-11 \\
WikiText-103 vs. MNLI & 0.7500 & 6.81e-12 \\
SST-2 vs. MNLI & 0.9839 & 3.94e-31 \\
\bottomrule
\end{tabular}
}
\end{table}

The three pairs involving WikiText-103 show tau values in [0.734, 0.786], while the three pairs among instruction-tuned and classification datasets show tau values in [0.940, 0.984]. Even the worst-case pair (WikiText-103 vs. SST-2) exceeds the gate threshold of 0.6 by a margin of 0.13.

ICC(3,k) sensitivity across epsilon values: 0.9862 (epsilon = 0.001), 0.9846 (epsilon = 0.01), 0.9878 (epsilon = 0.05), 0.9857 (epsilon = 0.1). All exceed 0.98. H-M1 gate result: \textbf{PASS}.

\begin{figure}[htbp]
\centering
\includegraphics[width=\columnwidth]{figures/sparsity_heatmap.png}
\caption{Sparsity heatmap}
\label{fig:sparsity_heatmap}
\end{figure}

\textit{Figure 3. Sparsity heatmap: 4 distributions by 32 layers at epsilon = 0.01. Rows are distributions; columns are layers.}

\begin{figure}[htbp]
\centering
\includegraphics[width=\columnwidth]{figures/pairwise_tau_matrix.png}
\caption{Pairwise tau matrix}
\label{fig:pairwise_tau_matrix}
\end{figure}

\textit{Figure 4. Pairwise Kendall's tau matrix for all four calibration distributions at epsilon = 0.01.}

\subsubsection{5.3 RQ3: Threshold Invariance (H-M2: PASS)}

All six cross-epsilon pairwise Kendall's tau values exceed 0.95 for epsilon values spanning a 100-fold range (0.001 to 0.1). The minimum is tau = 0.9597 (epsilon = 0.01 vs. epsilon = 0.1; p = 7.94e-27). The maximum is tau = 0.9960 (epsilon = 0.001 vs. epsilon = 0.01; p = 2.43e-34). CV exceeds 0.3 for all four epsilon values (pass rate 4/4).

\begin{table}[htbp]
\centering
\resizebox{\columnwidth}{!}{%
\begin{tabular}{lll}
\toprule
Pair & Kendall's tau & p-value \\
\midrule
epsilon = 0.001 vs. 0.01 & 0.9960 & 2.43e-34 \\
epsilon = 0.001 vs. 0.05 & 0.9677 & 4.47e-28 \\
epsilon = 0.001 vs. 0.1 & 0.9637 & 1.96e-27 \\
epsilon = 0.01 vs. 0.05 & 0.9718 & 9.26e-29 \\
epsilon = 0.01 vs. 0.1 & **0.9597** & 7.94e-27 \\
epsilon = 0.05 vs. 0.1 & 0.9798 & 2.82e-30 \\
\bottomrule
\end{tabular}
}
\end{table}

H-M2 gate result: \textbf{PASS}.

\begin{figure}[htbp]
\centering
\includegraphics[width=\columnwidth]{figures/cross_epsilon_tau_heatmap.png}
\caption{Cross-epsilon tau heatmap}
\label{fig:cross_epsilon_tau_heatmap}
\end{figure}

\textit{Figure 5. Cross-epsilon Kendall's tau matrix for all 6 pairs across epsilon in {0.001, 0.01, 0.05, 0.1}.}

\subsubsection{5.4 Gate Summary}

\begin{table*}[htbp]
\centering
\resizebox{\dimexpr 2\columnwidth+\columnsep\relax}{!}{%
\begin{tabular}{lllll}
\toprule
Hypothesis & Metric & Value & Threshold & Gate \\
\midrule
H-E1 & CV (epsilon = 0.01) & 0.544 & > 0.3 & PASS \\
H-E1 & tau\_calibration (Alpaca vs. WikiText) & 0.786 & >= 0.6 & PASS \\
H-M1 & ICC(3,k) & 0.9846 & > 0.75 & PASS \\
H-M1 & tau\_min (6 pairs) & 0.7339 & >= 0.6 & PASS \\
H-M2 & CV pass rate & 4/4 & >= 3/4 & PASS \\
H-M2 & Max adjacent-pair tau & 0.9960 & >= 0.7 & PASS \\
\bottomrule
\end{tabular}
}
\end{table*}

Prediction P1 (structural fingerprint exists and is stable): \textbf{SUPPORTED}. Predictions P2 (sparsity-rank correlation) and P3 (end-to-end allocation efficiency): \textbf{NOT TESTED} (H-M3 and H-M4 not executed).

